\documentclass[12pt]{article}
\usepackage[T1]{fontenc} 
\usepackage[utf8x]{inputenc} 
\usepackage{lmodern}
\usepackage [frenchb]{babel}
% Pour pouvoir utiliser 
\usepackage{ucs}

\usepackage{textcomp}
\usepackage{graphicx}
\usepackage{keystroke}
\usepackage{amssymb}
\usepackage{amsmath}
%\renewcommand{\thesection}{\arabic{section}} % numérotation des sections	
\usepackage[cc]{titlepic} %rajouter le logo dans la page de garde
\usepackage{url} % Pour avoir de belles url
\usepackage[a4paper]{geometry}
\usepackage[linktocpage]{hyperref}
% Pour pouvoir faire plusieurs colonnes
\usepackage {multicol}
%Pour faire plusieurs lignes
\usepackage{multirow}
\usepackage{slashbox}
% Pour mettre du code source
\usepackage {listings}
% Pour pouvoir passer en paysage
\usepackage{lscape}	

% Pour pouvoir faire plusieurs colonnes
\usepackage {multicol}

% POur crééer un index
\usepackage{makeidx}

%insertion de pdf
\usepackage{pdfpages}

%Utilisation d'acronyme
\usepackage{acronym}

\usepackage{fancyhdr}

\pagestyle{fancy}
\usepackage{lastpage} %numérotation type k/n
\renewcommand{\footrulewidth}{0.5pt} % trait horizontal
%\renewcommand{\headrulewidth}{1pt} % suppresion du trait horizontal dans l'entête
\lfoot{Analyse of a reseach article}
\cfoot{\thepage \hspace{0.15cm}sur  \pageref{LastPage} } % numérotation des pages
\rfoot{A.\textsc{M} \& M.\textsc{O}}
%\fancyhead[L]{\hspace{-1.5cm}\includegraphics[scale=0.105]{img/logouniv.eps}} %logo entête !
\fancyhead{\empty} %logo entête !


\fancypagestyle{plain}{% 1ères pages des chapitres:
   \fancyhead{} % supprime l’entete...
   \renewcommand{\headrulewidth}{0pt} % ...et le filet
}

\renewcommand{\thesection}{\arabic{section}}





\hypersetup{
  backref=true,
  %permet d'ajouter des liens dans...
  pagebackref=true,%...les bibliographies
  hyperindex=true, %ajoute des liens dans les index.
  colorlinks=true, %colorise les liens
  breaklinks=true, %permet le retour à la ligne dans les liens trop longs
  urlcolor= blue, %couleur des hyperliens
  citecolor=	cyan,
  bookmarks=true, %créé des signets pour Acrobat
  bookmarksopen=true,
  %si les signets Acrobat sont créés,
  %les afficher complètement.
  pdftitle={ScProjectArticle}, %informations apparaissant dans
  pdfauthor={AM and MO},
  %dans les informations du document
  pdfsubject={Doc}
  %sous Acrobat.
}

%\makeindex
%%%% debut macro pour enlever le nom chapitre %%%%
\makeatletter
\def\@makechapterhead#1{%
  \vspace*{50\p@}%
          {\parindent \z@ \raggedright \normalfont
            \interlinepenalty\@M
            \ifnum \c@secnumdepth >\m@ne
            \Huge\bfseries \thechapter\quad
            \fi
            \Huge \bfseries #1\par\nobreak
            \vskip 40\p@
}}

\def\@makeschapterhead#1{%
  \vspace*{50\p@}%
          {\parindent \z@ \raggedright
            \normalfont
            \interlinepenalty\@M
            \Huge \bfseries  #1\par\nobreak
            \vskip 40\p@
}}
\makeatother
%%%% fin macro %%%%
%Couverture 
\widowpenalty=10000
\clubpenalty=10000


\title{ {\huge \textsc{Senior capstone project}} \\  Master ALMA  \\ \textbf{Analyse of a research article} \\ \emph{Characterizing Software Architecture Changes : A Systematic Review}\\{\small Supervisor :  M.\textsc{Oussalah}}}
\author{A.\textsc{Marguerite} \& M.\textsc{Ouairy}}
\newcommand{\cf}{\emph{cf}.}


\begin{document}

\maketitle
\renewcommand{\labelitemi}{$\bullet$} 



\section{Introduction}\label{sec:intro}
This article was written by Byron J. \textsc{Williams} who is an assistant professor of the Department of Computer Science at the \href{http://www.msstate.edu/}{Mississippi State University} and Jeffrey C. \textsc{Carver} assistant professor too of the Department of Computer Science at the \href{http://www.ua.edu/}{University of Alabama}. The research focus area of this article are changes and their impact of the architecture of an information system. By this way, the authors \textbf{put forward} in the first part of the article a systematic literature review, that identifies and characterizes software changes (\cf{} section \ref{sec:pblm}). Then, in a second part, they introduce the \emph{Software Architecture Change Characterization Scheme} (SACCS). We will go into this contribution in greater depth in the section \ref{sec:res}. 


%% Cet article a été rédigé par M.\textsc{Butler} chercheur à l'\href{http://www.southampton.ac.uk/}{Université de Southampton} (UK) et K.\textsc{Damchoom} actuellement en Postdoc  à l'université de \href{http://www.southampton.ac.uk/}{Université de Southampton} (UK). Le travail proposé par ces deux chercheurs porte autour des méthodes de spécifications et de raisonnements formels. Plus particulièrement, ils utilisent la méthode \emph{Event-B}. Il s'agit d'une évolution de la méthode B introduisant la notion d'événements. Le remplacement des \og opérations \fg{} de la méthode B par des actions dotée1s de gardes donne la possibilité d'étudier des systèmes orientés plus \emph{hardware}. C'est la démarche choisie par les auteurs  dont la problématique sera développée dans la section \ref{sec:pblm}. 

\section{Research area}\label{sec:pblm}
Nowadays, developers know an increasing demands on software. These demands involve new software as well the upgrades of existing ones. Such operations lead developers to face up whith complexities and the risks of degrading the software architecture. A systematic review is a repeatable and formalized process in wich researcher systematically search a body of literature to document the state of knowledge on a particular subject. A such report, give to reseachers who use it a better \textbf{hindsight} than an ad-hoc review. \textsc{Williams} and \textsc{Carver} 

%% M.\textsc{Butler} et K.\textsc{Damchoom} proposent dans ce papier d'utiliser les méthodes décrites dans la section \ref{sec:intro} avec un exemple d'application concrète. Il s'agit de l'utilisation d'une mémoire flash, et de la gestion d'erreurs de ses opérations. Par exemple comment éviter la perte de données lors d'une coupure de courant, ou suite à un retrait manuel au cours d'une opération ? Suite à des éventuelles erreurs, quels mécanismes mettre en place pour garantir la consistance des données ? Les auteurs proposent d'étudier ce problème à partir d'une architecture flash (\href{http://www.onfi.org}{ONFI}) en mettant en oeuvre les outils proposés par la méthode Event-B avec l'aide de l'IDE (Rodin) associé.


\section{Main results}\label{sec:res}
%% Dans un premier temps, les auteurs présentent leur stratégie modélisation formelle du problème étudié. Ils utilisent notamment des raffinements. Deux types de raffinements sont distingués. Un raffinement vertical qui découpe le problème d'étude en 5 couches : de la première plus abstraite (\textit{Layer Abstraction}), à la cinquième moins abstraite (\textit{Fourth refinement}). Plus on \og descend \fg{} dans les couches concrètes, plus on se rapproche du modèles réel en rajoutant de nouvelles contraintes. Un raffinement horizontal consiste notamment à l'introduction de nouvelles étapes dans les opérations étudiées. Par exemple une opération d'écriture \verb+write+ est décomposée en trois sous étapes : \verb+begin+, \verb+step+ et \verb+end+. Ces décompositions en sous actions donnent de nouvelles possibilité pour proposer des preuves du bon déroulement du protocole étudié.
%% \paragraph{}
%% Dans un second temps, M.\textsc{Butler} et K.\textsc{Damchoom} présentent les différents algorithmes en Event-B en appliquant la stratégie précédemment décrite. Dans ce cours résumé, nous nous contenterons de présenter les idées clefs de certains d'entre eux. La consistance des données est garantie par un mécanisme utilisant des mémoires/étapes \og{} tampons\fg{}. À titre d'exemple, l'écriture d'un fichier (décomposé en bloques de datas) se fait dans un premier temps dans une variable tampon ($f\_compt\_temp$). Si chaque étape/événement (dotés de garde) de cette procédure se déroule sans erreurs, alors la variable tampon devient le résultat de la copie ($f\_compt \Rightarrow f\_content$). Sinon on passe dans un état d'erreur empêchant l'écriture de données corrompues.
%% \paragraph{}
%% Un autre exemple de problématiques à laquelle ont répondu les auteur, est la gestion des bloques mémoire de la mémoire flash. En effet les nombreuses opérations sur ces derniers entraînent l'isolement de bloques \og corrompus \fg{} (non vide, mais ne contenant pas de données valides). La stratégie mise en place est de \og déplacer \fg{} toutes les pages de donnée mémoire valides régulièrement à un autre endroit de la mémoire. Une fois ce travail achevé tous les autres bloques sont considérés comment vides et peuvent être écrasés.
\clearpage

\section{Issue remaining}
%% Bien que satisfaits par les résultats des preuves automatiques (aux alentours de 93\%), les auteurs restent face à certains problèmes qui une fois résolus permettraient une augmentation de ce taux de réussite. La nécéssité de reprouver certaines preuves à chaque occurence d'un changement de de modèle contribue à diminuer ce taux de réussite. Les auteurs expliquent que l'ajout de certaines fonctionalités manquante de l'outil Rodin pourrait résoudre certains de ces problèmes. De la même manière l'outil ne propose pas directement les découpages nécéssaires au rafinements. Ce sont donc les auteurs quis opèrent manuellement ce découpage. Ils émettent alors l'idée de la créaton d'un outil dédié à cette décomposition. 


\section{Discussion an re-use}
%% Cet article décrit de manière pédagogique et détaillée des algorithmes en Event-B fournissant des résolutions automatiques de preuves de l'utlisation sure d'une mémoire flasj avec de très bons résultats. Les différents raffinements proposés par M.\textsc{Butler} et K.\textsc{Damchoom} entraînent le découpages du problème étudié en de multiples sous problèmes plus \og petits \fg{}. Ainsi les résolutions de ces sous problèmes apparaissent simples et efficaces. Faces à ces bons résultats, on peut envisager la réutilisation de ces découpages \og multi-niveaux \fg{} pour d'autres problématiques.
%% De plus cet article met en exergue l'efficacité de la méthode Event-B dans le découpage d'événements atomiques. Des réutilisations de ces méthodes pour les problèmes de concurrences ou de tolérances aux pannes (protocoles de communications, systèmes embarqués temps réels~\dots) peuvent être envisagées.

\end{document}

%%  LocalWords:  Butler Southampton UK Damchoom Postdoc Event-B ONFI
%%  LocalWords:  l'IDE Fourth refinement skip end fail compt write
%%  LocalWords:  begin step multi temp SACCS
